\chapter{Limitations, Scope for Improvement and Future Work}
\label{ch:limitations}

\section{Limitations}

\paragraph{Theoretical assumptions.}
\begin{itemize}
\item T1 assumes infinite plates and a slowly varying planform. It is invalid near S5 and in the outlet.
\item T2 assumes locally fully developed flow in each rectangular section.
\item T3 is depth-averaged and does not resolve the side-wall layer or three-dimensional effects at the corners.
\item All three are Stokes-flow models. This is justified by the CFD, which shows a non-linearity of at most \SI{0.7}{\percent}.
\end{itemize}

\paragraph{Idealised boundary conditions.} The inflow enters as a uniform velocity at the face of a \SI{2}{mm} lead-in. The punched ports, tubing and bubble trap of the real device were not modelled. The outlet is at a fixed pressure. These choices do not affect the hyperbolic section, but they do affect the absolute pressures.

\paragraph{Numerical discretisation and mesh resolution.} Only one mesh was run. About nine cells span the gap and the wall nodes are \SIrange{0.08}{0.10}{mm} apart near the side walls. The uniform \SI{-5}{\percent} offset of shear and pressure gradient is most probably discretisation error, but this has not been demonstrated. The mesh type and quality metrics were not documented in the exports.

\paragraph{Missing data and convergence evidence.} The channel height and the lead-in were inferred rather than read from the set-up. The velocity was exported without the $z$ coordinate. Neither the residual history nor the mass balance was exported.

\paragraph{Reference data.} The published results are available only as plots and ranges. The paper's height, length and uniformity criterion are partly inconsistent, so the comparison with it is semi-quantitative.

\paragraph{Oxygen model.}
\begin{itemize}
\item The uptake parameters ($\qmax$, $\Km$) are illustrative, not measured for Caco-2 or A549 cells, and the conclusions depend on $\qmax$ above all.
\item The model is steady. It ignores re-oxygenation transients after medium changes and the hours needed to equilibrate the PDMS.
\item The PDMS is represented by a steady Robin coefficient. Partitioning, diffusion inside the PDMS, water immersion and permeability loss by protein adsorption are not modelled.
\item The reduced-order model is uniform across the width. It neglects the slower side-wall layers, where depletion with impermeable walls will be stronger.
\item The monolayer is a surface flux. Cell height, its effect on the flow and the variation of uptake with shear or with cell differentiation are ignored.
\item There is no experimental validation of any oxygen prediction.
\end{itemize}

\section{Future work}

Each item below addresses one of the limitations above.

\begin{enumerate}
\item \textbf{Mesh-independence study (highest priority).} Run the two-factor study of \cref{sec:num_gci} at \SI{1}{\uLmin} ($N_z=9$, 18, 36; $\Delta s=100$, 50, \SI{25}{\micro\metre}) and report $\mathrm{GCI}_{\mathrm{fine}}$, the observed order and the extrapolated values. If the shear rises towards T2 with $N_z$, the \SI{-5}{\percent} offset is confirmed as discretisation error and the absolute CFD values can be corrected.
\item \textbf{Complete exports.} Report the inlet/outlet mass flux, the residual history and the domain extents. Export $u_x(z)$ along lines at S1--S5, so the parabola and a profile-fit wall shear can be tested, and export the floor ($z=0$) wall shear alone along fine lines across the width.
\item \textbf{Oxygen CFD.} Run the prepared Fluent case: one user-defined scalar (UDS) on the frozen flow, with the Michaelis--Menten floor flux and the PDMS Robin flux implemented as user-defined functions (UDFs). Start with the mandatory sign and mass-balance test (zero-order uptake with impermeable walls must satisfy $Q(\Csat-\Cb_{\mathrm{out}})=\qmax A_{\mathrm{cells}}$). Then run IMP, PDMS and SAT at the three flow rates, and compare with \cref{tab:o2_results}. The CFD will quantify the extra depletion in the side-wall layers and near the outlet corner, which the reduced-order model cannot capture.
\item \textbf{Transient oxygen transport.} Simulate re-oxygenation after a medium change or a change of flow, over 1--\SI{2}{h}.
\item \textbf{Conjugate PDMS transport.} Model oxygen transport inside the lid and side walls (solved in partial pressure, to handle partitioning). Vary the lid thickness, permeability ageing and water immersion.
\item \textbf{Measured kinetics.} Measure the oxygen consumption rate and the confluent density of the cell lines, and replace the illustrative $\qmax$ and $\Km$. Consider uptake that depends on the shear or on the differentiation state of the cells.
\item \textbf{Experimental validation.} Measure the pressure drop (flow), and the oxygen concentration at the floor using an oxygen-sensing film (for example a PtTFPP optode) along the channel at 0.22 and \SI{1}{\uLmin}.
\item \textbf{Parametric design.} Use the closed forms to choose $h$, $w_1$, $L$ and $Q$ that meet a target shear range while keeping $\Da_{\mathrm{conv}}<1$ or $\Phi_{\mathrm{PDMS}}\gg1$. Check selected designs with CFD.
\end{enumerate}
